\section{Komplexe Analysis}

\subsection{Die komplex konjugierte Geschwindigkeit $\overline{v(z)}$}

Um den Ausdruck~\eqref{joukowski:auftrieb:auftrieb} weiter aufzulösen, muss die Funktion $v(z)$ genauer untersucht werden.
Einige Eigenschaften von $v(z)$ helfen uns dabei einen passenden Ausdruck dafür aufzustellen.
Wir haben vorausgesetz, dass im Unendlichen die Geschwindigkeit $v(z) = v_\infty \in \mathbb{R}$ ergibt.
Zudem muss es laut Beispiel~\ref{buch:integration:wege:bsp:kreisintegral} eine einfache Polstelle im inneren des Flügelprofils haben,
da sonst das Kontourintegral konstant $0$ wäre,
was darauf hindeuten würde, dass keinen Auftrieb erzeilt werden kann.
Diese Eigenschaften der Geschwindigkeit $v(z)$ führen dazu, dass die komplex konjugierte Geschwidigkeit $\overline{v(z)}$ als Laurent-Reihe der Form

\[
	\overline{v(z)} 
	= 
	a_0 + \frac{a_1}{z} + \frac{a_2}{z^2} + \frac{a_3}{z^3} + \dots 
	= 
	\sum_{k=0}^{\infty} \frac{a_k}{z^k}
\]
geschrieben werden kann. Im Allgemeinen sind die Koeffizienten $a_k \in \mathbb{C}$, in dieser Anwendung jedoch beschränken sich die Koeffizienten auf $a_k \in \mathbb{R}$.
Da die Geschwindigkeit $v(z)$ überall ausser im Punkt $z_0 = 0$ definiert ist, muss auch der Reihenausdruck diese Eigenschaft haben.
Dies wird zum Beispiel von einer Taylor-Reihe nicht erfüllt, weshalb die Laurent-Reihe gewählt wurde.

Der Koeffizient $a_0$ entspricht direkt der Geschwindigkeit im Unendlichkeit, weshalb er besser als $v_\infty$ geschrieben werden soll.
Für die komplex konjugierte Geschwindigkeit $\overline{v(z)}$ ergibt sich somit die besser formulierte Reihenentwicklung

\begin{equation}
	\overline{v(z)}
	=
	v_\infty + \sum_{k=1}^{\infty} \frac{a_k}{z^k}.
	\label{joukowski:komplexe-analysis:laurent-reihe}
\end{equation}

% Zusammenhang von |v(z)|^2 und v(z)*^2?

\subsection{Komplexe Kraft}

Nun kann die Gleichung~\eqref{joukowski:komplexe-analysis:laurent-reihe} in den Ausdruck~\eqref{joukowski:auftrieb:auftrieb} eingesetz werden
und in ein Integral mit allen Termen die mit $v_\infty$ multipliziert werden, 
und ein Integral mit allen anderen Termen umformen, wobei aus

\[
	\begin{aligned}
		F'
		&= 
		-i\frac{\rho}{2} \cdot 
		\displaystyle\oint_{\gamma} v_\infty^2 + 2 v_\infty \sum_{k=1}^{\infty} \frac{a_k}{z^k} + \biggl(\sum_{k=1}^{\infty} \frac{a_k}{z^k}\biggr)^2
		dz \\
		\quad
		&= 
		-i\frac{\rho}{2} \cdot  
		\displaystyle\oint_{\gamma} 
		2 v_\infty \biggl(\frac{v_\infty}{2} + \sum_{k=1}^{\infty} \frac{a_k}{z^k} \biggr)
		dz
		-
		i\frac{\rho}{2} \cdot 
		\displaystyle\oint_{\gamma} 
		\biggl(\sum_{k=1}^{\infty} \frac{a_k}{z^k}\biggr)^2
		dz
	\end{aligned}
\]
wird. 
Beim ersten Ausdruck kann der Term $\frac{v_\infty}{2}$ in $v_\infty - \frac{v_\infty}{2}$ umwandeln. 
Somit lässt sich $-\frac{v_\infty}{2}$ ins zweite Integral verschieben.
Jetzt kann 

\[
	F'
	= 
	-i \rho v_\infty \cdot
	\displaystyle\oint_{\gamma} 
	\biggl(v_\infty + \sum_{k=1}^{\infty} \frac{a_k}{z^k}\biggr)
	dz
	-
	i\frac{\rho}{2} \cdot 
	\displaystyle\oint_{\gamma} 
	\biggl(\sum_{k=1}^{\infty} \frac{a_k}{z^k}\biggr)^2 - \frac{v_\infty^2}{2}
	dz
\]
mit der Gleichung ~\eqref{joukowski:komplexe-analysis:laurent-reihe} in
\begin{equation}
	F'
	= 
	-i \rho v_\infty \cdot
	\displaystyle\oint_{\gamma} \overline{v(z)}
	dz
	-
	i\frac{\rho}{2} \cdot 
	\displaystyle\oint_{\gamma} 
	\biggl(\sum_{k=1}^{\infty} \frac{a_k}{z^k}\biggr)^2 - \frac{v_\infty^2}{2}
	dz
	\label{joukowski:komplexe-analysis:kraft-mit-zwei-termen}
\end{equation}
geschrieben werden.
Da $\frac{v_\infty}{2}$ konstant ist, verschwindet es beim Integrieren um eine geschlossene Kurve.
Somit können wir für den Ausdruck des ersten Integrals aus der Gleichung~\eqref{joukowski:komplexe-analysis:kraft-mit-zwei-termen} die \emph{Zirkulation} definieren.

\begin{definition}[Zirkulation]
	\label{joukowski:komplexe-analysis:def:zirkulation}
	Die Zirkulation $\Gamma$ ist das Mass für die um ein Profil drehende Strömung und ist das Ergebnis des Linienintegrals
	\[
		\Gamma
		=
		\displaystyle\oint_{\gamma} \overline{v(z)}
		dz.
	\]
\end{definition}

Dem zweiten Integral der Gleichtung~\eqref{joukowski:komplexe-analysis:kraft-mit-zwei-termen} geben wir den Namen $\Lambda$.
somit kann die Kraft $F'$ geschrieben werden als

\begin{equation}
	F'
	=
	-i \rho v_\infty \Gamma - i \frac{\rho}{2} \Lambda.
	\label{joukowski:komplexe-analysis:kraft-mit-lambda}
\end{equation}

Das Problem können wir in zwei Teile aufteilen:
Einerseits die Auswertung von $\Lambda$,
andererseits die Berechnung der Zirkulation.

\subsubsection{Berechnung von $\Lambda$}
\label{joukowski:komplexe-analysis:berechnung-von-lambda}

Für die Auswertung von $\Lambda$
aus der Gleichung~\eqref{joukowski:komplexe-analysis:kraft-mit-lambda} können wir ausnutzen,
dass nicht um die Kontour des Flügelprofils direkt integriert werden muss,
sondern dass jeder Weg, welcher das Flügelprofil komplett umschliesst, dasselbe Resultat gibt.
Wir wählen als Weg den Kreis mit Zentrum bei $z_0 = 0$ und mit Radius $r$.
Dafür muss jetzt die Definition~\ref{buch:integration:definition:definition:wegintegral} des Wegintegrals verwendet werden.
Wobei wir die Funktion $f(z) = \biggl(\displaystyle\sum_{k=1}^{\infty} \frac{a_k}{z^k}\biggr)^2$ wählen.
Dabei parametrisieren wir $z = \gamma(t) = r e^{i t}$ und stellen das Integral 
\begin{equation}
	\begin{aligned}
		\Lambda
		&=
		\displaystyle\int_{0}^{2\pi} 
		\biggl(\sum_{k=1}^{\infty} \frac{a_k}{r^k e^{k \cdot i t}}\biggr)^2
		i r e^{i t}
		dt \\
		\quad
		&=
		i \displaystyle\int_{0}^{2\pi}
		\bigg(\frac{a_1^2}{r^2 e^{2i t}} + \frac{a_1 a_2}{r^3 e^{3i t}} + \frac{a_2^2}{r^4 e^{4i t}} + \dots \bigg) 
		r e^{i t}
		dt \\
		\quad
		&=
		i \displaystyle\int_{0}^{2\pi}
		\bigg(\frac{a_1^2}{r e^{i t}} + \frac{a_1 a_2}{r^2 e^{2i t}} + \frac{a_2^2}{r^3 e^{3i t}} + \dots \bigg) 
		dt
	\end{aligned}
	\label{joukowski:komplexe-analysis:parametrisiertes-lambda}
\end{equation}
auf.
Da die Auswertung des Integrals nicht auf den Radius des Kreises darauf ankommt, sofern er das ganze Fügelprofil umschliesst,
können wir den Grenzübergang des Radiuses $\displaystyle\lim_{r\to \infty}$ durchführen.
Weil jeder Term der Gleichung~\eqref{joukowski:komplexe-analysis:parametrisiertes-lambda} 
ein $r$ der Ordnung $1$ bis $n$ im Nenner hat,
gehen alle gegen $0$.

Somit muss $\Lambda = 0$ sein.
Der Term~\eqref{joukowski:komplexe-analysis:kraft-mit-lambda} verkleinert sich zur Form $F' = -i \rho v_\infty \Gamma$, 
welche als Satz von Kutta-Joukowski bekannt ist.

\begin{satz}[Satz von Kutta-Joukowski]
	Der Satz von Kutta-Joukowski beschreibt die Proportionalität zwischen der komplexen Kraft und der Zirkulation,
	welche sich mit
	\[
		F'
		=
		-i \rho v_\infty \Gamma
	\]
	berechnen lässt.
\end{satz}

\subsubsection{Berechnung der Zirkulation $\Gamma$}
Für die Auswertung von $\Gamma$ kann sehr ähnlich vorgegangen werden, wie bei der Berechnung von $\Lambda$ in Abschnitt~\ref{joukowski:komplexe-analysis:berechnung-von-lambda}.
Auch da wählen wir den Kreis mit Zentrum bei $z_0 = 0$ und mit Radius $r$ und die gleiche parametrisierung von $z$.
Für die Funtion $f(z) = \overline{v(z)}$ setzen wir die Reihenentwicklung aus der Gleichung~\eqref{joukowski:komplexe-analysis:laurent-reihe} ein.
Somit wird
\[
	\begin{aligned}
		\Gamma
		&=
		\displaystyle\int_{0}^{2\pi}
		\biggl(v_\infty + \sum_{k=1}^{\infty} \frac{a_k}{r^k e^{k \cdot i t}}\biggr)
		i r e^{i t}
		dt \\
		\quad
		&=
		\displaystyle\int_{0}^{2\pi}
		\biggl(v_\infty + \frac{a_1}{r e^{i t} + \frac{a_2}{r^2 e^{2 i t}} + \frac{a_3}{r^3 e^{3 i t}} + \dots \biggr)
		i r e^{i t}
		dt \\
		\quad
		&=
		i\displaystyle\int_{0}^{2\pi}
		\biggl(v_\infty r e^{i t} + a_1 + + \frac{a_2}{r e^{i t} + \frac{a_3}{r^2 e^{2 i t}} + \dots \biggr)
		dt
	\end{aligned}
\]